\section{Introduction}

Researchers who study large collections of text rarely read all of them. An annotation team codes a few thousand posts out of millions, a large language model summarizes a few hundred documents because the whole corpus does not fit in its context window, and a qualitative researcher keeps reading until further documents stop adding codes. Each of these practices assumes that the documents in hand stand for the documents left unread. Yet the number of documents is usually set by budget or by feel. If we read 15 posts from a corpus of 10{,}000, a number a person can read in a few minutes, have we understood the corpus? If not, would 100 posts be enough, or 1{,}000?

Survey research has a settled answer to the same question for a human population. The researcher states how much people differ on the quantity of interest (the population variance $\sigma^{2}$) and how large an error is tolerable ($e$), and the formula $n = z^{2}\sigma^{2}/e^{2}$ gives the number of people to interview \citep{cochran1977, kish1965}. Every estimate from the sample then carries a margin of error. A sample of a corpus comes with neither a sample size formula nor a margin of error. This paper builds both for the question a reader of a sample usually has in mind: how much of the corpus has a document like it among the documents we read?

The answer depends on what counts as alike. Two posts can be about the same topic, such as the 2020 United States Census, while making different claims, and two other posts can repeat the same claim in different words. We place each document in an embedding space, where a text encoder puts documents with similar meaning close together, and we measure how alike two documents are by the cosine distance between them. The user then chooses a distance $r$ below which two documents count as the same content. To give $r$ a concrete meaning, we asked two large language models, Qwen3.8 and Gemma~4, to judge about 4{,}000 pairs of documents at a range of distances. For each corpus we found the distance below which most pairs are about the same specific claim, question, event or story, and the larger distance below which most pairs share at least a topic (Section~\ref{sec:rcal}).

At these two distances the question of the 15 posts has a clear answer, and it differs from corpus to corpus (Table~\ref{tab:reader}). Take the 17{,}202 tweets about the 2020 Census that contain the word ``illegal''. In a random sample of 15 of them, 98\% of the corpus has a tweet on the same topic, so a reader of those 15 tweets has seen nearly every topic the corpus discusses. Only 38\% of the corpus, however, has a tweet in the sample that makes the same specific point. Now take 10{,}000 personal stories from the TIFU forum on Reddit, in which users tell how they messed something up. In a sample of 15 stories, only 1.5\% of the corpus has a story on the same topic. Even in a sample of 1{,}000 stories, only 10\% of the corpus has a story about the same event. Whether a sample is enough therefore depends on the corpus and on what the reader needs to know.

The share of a corpus that has no document like it in the sample is called the missing mass, and statistics has a classical way to estimate it without looking at the unread documents. During the Second World War, I.~J. Good and Alan Turing noticed that a sample reveals how much it is missing through the items it contains only once \citep{good1953}. If most words in a sample of text occur once, many more words remain unseen. If every word in the sample has occurred several times, few words remain. \citet{maurer2022missing} carried this idea from categories to distances. His estimator $G$, which we call the isolation rate, is the share of documents in the sample that have no other sampled document within $r$, that is, the share of the documents we read that were unlike anything else we read. For example, in a sample of 15 of the Census tweets with ``illegal'', 64\% of the tweets have no close match among the other 14. This is the estimate. The truth, which we can compute because we have the whole corpus, is that 62\% of the corpus has no close match in the sample. A reader can therefore learn how much a sample has missed from the sample alone. Maurer's paper is theoretical. It proves how far $G$ can be from the missing mass, and it leaves open how $G$ falls as the sample grows, which is what a researcher needs in order to plan a sample. To our knowledge the estimator has not yet been used on text.

How $G$ falls with the sample size follows from a classical result about the distances between random points \citep{levina2005, facco2017twonn, penrose2003rgg}. Suppose the documents of a corpus differ from one another along $d$ independent directions of meaning, where $d$ is called the intrinsic dimension of the embedded corpus and can be estimated from the distances between documents. Then the chance that a document has no neighbour within $r$ among $n$ sampled documents depends on $n$ and $r$ only through the product $n r^{d}$. In other words, a larger sample and a looser definition of the same content can stand in for each other, at a rate set by $d$. The product has a direct meaning for sampling. When the definition of the same content becomes twice as strict, that is, when $r$ is halved, the number of documents close enough to count as a match shrinks $2^{d}$ times, so the sample must grow $2^{d}$ times to cover the corpus as well as before. For the medical questions of the AskDocs forum and the TIFU stories, $d$ is about 13, so a definition twice as strict requires about 8{,}000 times as many documents. The intrinsic dimension thus does for a corpus what the variance does for a survey population. It is the property that sets how large a sample must be.

We test this law on seventeen pairs of corpus and encoder, which include medical questions, personal stories, election discussion, tweets about the 2020 Census, posts from Truth Social and Bluesky, and synthetic data with a known answer. The law holds on most of these corpora, it needs one correction for finite corpora, and it fails on a recognizable kind of corpus.

\begin{enumerate}
\item On ten of the seventeen pairs, $G$ depends on the sample size and the distance only through $n r^{d}$. The law has two constants besides $d$, and a fit of these constants on samples of up to 1{,}000 documents predicts $G$ for samples seven times larger. Solved for $n$, the law gives the sample size that brings $G$ down to a target, such as 20\% of the sampled documents isolated. We tested this formula in 83 cases, one for each combination of corpus, target and distance at which the true answer lay beyond 1{,}000 documents. In 77 of the 83 cases, the sample that actually reached the target held between two thirds of the formula's number and one and a half times that number. A formula answer of 1{,}500 documents thus meant an actual need between 1{,}000 and 2{,}250. The standard tool from ecology for the same task, the extrapolation of \citet{chao2012coverage}, predicts $G$ for a sample twice as large with an average error two to nine times that of the law.
\item For a corpus of $N$ documents, the missing mass is $G\,(1 - n/N)$. The correction is needed because $G$ describes the documents a sample has not reached, and the $n$ sampled documents are reached by definition. Survey sampling knows the same factor as the finite population correction. With it, the estimate is on average one percentage point from the true missing mass.
\item The law fails on the other seven pairs. Five are political corpora, among them a forum in which short slogans repeated many times sit beside longer discussion. One is a second encoding of that forum, and one is a synthetic mixture of topics. A test run on a pilot sample of 2{,}000 documents separates the seven failing pairs from the ten passing pairs before the full sample is drawn.
\end{enumerate}

Taken together, these results let a researcher size a sample before reading it. We set a modest target, that half of the corpus should have a document with the same specific content in the sample. On corpora that repeat a few claims, such as Census tweets selected by a keyword or a forum about one election, a few hundred posts reach this target. On forum questions, personal stories and the Truth Social posts, the sample must hold more than 40\% of the corpus. The reference point is a corpus in which no two documents share content. There the sample matches only itself, so reaching the target takes exactly half of the corpus, and the forum questions, stories and Truth Social posts come close to this limit. Two pilots of different sizes serve two purposes. A pilot of 100 documents already tells which kind of corpus a researcher has, since in it no document has a match on the corpora that come close to the limit. A pilot of 2{,}000 documents is needed to fit the law and predict the required sample.
